	%  The AAU Poster Theme.
%  2013-05-08 v. 1.1.0
%  Copyright 2013 by Jesper Kjær Nielsen <jkn@es.aau.dk>
%
%  This is free software: you can redistribute it and/or modify
%  it under the terms of the GNU General Public License as published by
%  the Free Software Foundation, either version 3 of the License, or
%  (at your option) any later version.
%
%  This is distributed in the hope that it will be useful,
%  but WITHOUT ANY WARRANTY; without even the implied warranty of
%  MERCHANTABILITY or FITNESS FOR A PARTICULAR PURPOSE.  See the
%  GNU General Public License for more details.
%
%  You can find the GNU General Public License at <http://www.gnu.org/licenses/>. RATIO: 1,1165048544


% For UCSD EXPO
\documentclass[paperwidth=36in,paperheight=27in,landscape,11pt]{preambles/baposter}

% For IV 2015
%\documentclass[paperwidth=90cm,paperheight=130cm,landscape,11pt]{preambles/baposter}




%%%%%%%%%%%%%%%%%%%%%%%%%%%%%%%%%%%%%%%%%%%%%%%%
% Language, Encoding and Fonts
% http://en.wikibooks.org/wiki/LaTeX/Internationalization
%%%%%%%%%%%%%%%%%%%%%%%%%%%%%%%%%%%%%%%%%%%%%%%%
% Select encoding of your inputs. Depends on
% your operating system and its default input
% encoding. Typically, you should use
%   Linux  : utf8 (most modern Linux distributions)
%            latin1 
%   Windows: ansinew
%            latin1 (works in most cases)
%   Mac    : applemac
% Notice that you can manually change the input
% encoding of your files by selecting "save as"
% an select the desired input encoding. 
\usepackage[utf8]{inputenc}
% Make latex understand and use the typographic
% rules of the language used in the document.
\usepackage[english]{babel}
% Use the vector font Latin Modern which is going
% to be the default font in latex in the future.
\usepackage{helvet}
% Change the default font family from roman to sans serif
\renewcommand{\familydefault}{\sfdefault} % for text
\usepackage[helvet]{sfmath} % for math
% Choose the font encoding
\usepackage[T1]{fontenc}
%\usepackage[demo]{graphicx}
\usepackage{subfig}
\usepackage{multicol}
\usepackage{multirow}
\usepackage{blindtext}
\usepackage{booktabs}
\usepackage{multirow}
\usepackage{graphics} 
\usepackage{amssymb}
\usepackage{lipsum}
\usepackage{blindtext}
\usepackage{gensymb}
\usepackage{amsmath}
\usepackage{float}
%\usepackage{float}

%%%%%%%%%%%%%%%%%%%%%%%%%%%%%%%%%%%%%%%%%%%%%%%%
% Graphics and Tables
% http://en.wikibooks.org/wiki/LaTeX/Importing_Graphics
% http://en.wikibooks.org/wiki/LaTeX/Tables
% http://pgfplots.sourceforge.net/
%%%%%%%%%%%%%%%%%%%%%%%%%%%%%%%%%%%%%%%%%%%%%%%%
% You cannot use floats in the baposter theme.
% We therefore load the caption package which provides
% the command \captionof
% Set up how figure and table captions are displayed
\usepackage{caption}
\captionsetup{
justification=centering,
  font=scriptsize,% set font size to footnotesize
  labelfont=bf % bold label (e.g., Figure 3.2) font
}
\captionsetup[subfloat]{
	labelfont={scriptsize}
	,textfont={scriptsize}}
% Make the standard latex tables look so much better
\usepackage{array,booktabs}
\usepackage{multicol}
\usepackage{txfonts}


% For creating beautiful plots
\usepackage{pgfplots}
\usepackage{multirow}
\usepackage{rotating}
\usepackage{bchart}
\usepackage{newfloat}
\usepackage{wrapfig,lipsum,booktabs}
\DeclareFloatingEnvironment[
  fileext=loc
]{chart}
%%%%%%%%%%%%%%%%%%%%%%%%%%%%%%%%%%%%%%%%%%%%%%%%
% Mathematics
% http://en.wikibooks.org/wiki/LaTeX/Mathematics
%%%%%%%%%%%%%%%%%%%%%%%%%%%%%%%%%%%%%%%%%%%%%%%%
% Defines new environments such as equation,
% align and split 
\usepackage{amsmath}
% Adds new math symbols
\usepackage{amssymb}

%%%%%%%%%%%%%%%%%%%%%%%%%%%%%%%%%%%%%%%%%%%%%%%%
% Colours
% http://en.wikibooks.org/wiki/LaTeX/Colors
%%%%%%%%%%%%%%%%%%%%%%%%%%%%%%%%%%%%%%%%%%%%%%%%
\selectcolormodel{RGB}
% define the three aau colors
\definecolor{aaublue1}{RGB}{33,26,82}% dark blue
\definecolor{aaublue2}{RGB}{113,109,143} % light blue
\definecolor{aaublue3}{RGB}{194,193,204} % lighter blue
\definecolor{black}{RGB}{0,0,0}
\definecolor{white}{RGB}{255,255,255}

%%%%%%%%%%%%%%%%%%%%%%%%%%%%%%%%%%%%%%%%%%%%%%%%
% Lists
% http://en.wikibooks.org/wiki/LaTeX/List_Structures
%%%%%%%%%%%%%%%%%%%%%%%%%%%%%%%%%%%%%%%%%%%%%%%%
% Easier configuration of lists
\usepackage{enumitem}
%configure itemize
\setlist{%
  topsep=0pt,% set space before and after list
  noitemsep,% remove space between items
  labelindent=\parindent,% set the label indentation to the paragraph indentation
  leftmargin=*,% remove the left margin
  font=\color{aaublue1}\normalfont, %set the colour of all bullets, numbers and descriptions to aaublue1
}
% use set<itemize,enumerate,description> if you have an older latex distribution
\setitemize[1]{label={\raise1.25pt\hbox{$\blacktriangleright$}}}
\setitemize[2]{label={\scriptsize\raise1.25pt\hbox{$\blacktriangleright$}}}
\setitemize[3]{label={\raise1.25pt\hbox{$\star$}}}
\setitemize[4]{label={-}}
%\setenumerate[1]{label={\theenumi.}}
%\setenumerate[2]{label={(\theenumii)}}
%\setenumerate[3]{label={\theenumiii.}}
%\setenumerate[4]{label={\theenumiv.}}
%\setdescription{font=\color{aaublue1}\normalfont\bfseries}

% use setlist[<itemize,enumerate,description>,<level>] if you have a newer latex distribution
%\setlist[itemize,1]{label={\raise1.25pt\hbox{$\blacktriangleright$}}}
%\setlist[itemize,2]{label={\scriptsize\raise1.25pt\hbox{$\blacktriangleright$}}}
%\setlist[itemize,3]{label={\raise1.25pt\hbox{$\star$}}}
%\setlist[itemize,4]{label={-}}
%\setlist[enumerate,1]{label={\theenumi.}}
%\setlist[enumerate,2]{label={(\theenumii)}}
%\setlist[enumerate,3]{label={\theenumiii.}}
%\setlist[enumerate,4]{label={\theenumiv.}}
%\setlist[description]{font=\color{aaublue1}\normalfont\bfseries}

%%%%%%%%%%%%%%%%%%%%%%%%%%%%%%%%%%%%%%%%%%%%%%%%
% Misc
%%%%%%%%%%%%%%%%%%%%%%%%%%%%%%%%%%%%%%%%%%%%%%%%
% change/remove some names
\addto{\captionsenglish}{
  %remove the title of the bibliograhpy
  \renewcommand{\refname}{\vspace{-0.7em}}
  %change Figure to Fig. in figure captions
  \renewcommand{\figurename}{Figure}
}
% create links
\usepackage{url}
%note that the hyperref package is currently incompatible with the baposter class

%%%%%%%%%%%%%%%%%%%%%%%%%%%%%%%%%%%%%%%%%%%%%%%%
% Macros
%%%%%%%%%%%%%%%%%%%%%%%%%%%%%%%%%%%%%%%%%%%%%%%%
\newcommand{\alert}[1]{{\color{aaublue1}#1}}

%%%%%%%%%%%%%%%%%%%%%%%%%%%%%%%%%%%%%%%%%%%%%%%%
% Document Start 
%%%%%%%%%%%%%%%%%%%%%%%%%%%%%%%%%%%%%%%%%%%%%%%%
\begin{document}
\fontsize{12pt}{12pt}\selectfont
%%%%%%%%%%%%%%%%%%%%%%%%%%%%%%%%%%%%%%%%%%%%%%%%
% Some changes that cannot be made in the preamble
%%%%%%%%%%%%%%%%%%%%%%%%%%%%%%%%%%%%%%%%%%%%%%%%
% set the background of the poster
\background{
  \begin{tikzpicture}[remember picture,overlay]%
    %the poster background color
    \fill[fill=white] (current page.north west) rectangle (current page.south east);
    %the header
    \fill [fill=aaublue1] (current page.north west) rectangle ([yshift=-\headerheight] current page.north east);
  \end{tikzpicture}
}
% if you want to reduce the space before and after equations, use and adjust
% the following lines
%\addtolength{\abovedisplayskip}{-2mm}
%\addtolength{\belowdisplayskip}{-2mm}

%%%%%%%%%%%%%%%%%%%%%%%%%%%%%%%%%%%%%%%%%%%%%%%%
% General poster setup
%%%%%%%%%%%%%%%%%%%%%%%%%%%%%%%%%%%%%%%%%%%%%%%%
\begin{poster}{
  %general options for the poster
  grid=false,
  columns=3,
%  colspacing=4.2mm,
  headerheight=0.13\textheight,
  background=user,
%  bgColorOne=red!42, %is used when background != user and none
%  bgColortwo=green!42, %is used when background is shaded
  eyecatcher=true,
  %posterbox options
  headerborder=closed,
  borderColor=black,
  headershape=roundedright,
  headershade=plain,
  headerColorOne=aaublue2,
%  headerColortwo=yellow!42, %is used when the header background is shaded
  textborder=rectangle,
  boxshade=plain,
  %boxColorOne=black!10,
  boxColorOne=white,
%  boxColorTwo=cyan!42,%is used when the text background is shaded
  headerFontColor=white,
  headerfont=\Large\sf\bf,
  linewidth=1pt
}
%the Eye Catcher (the logo on the left)
{
  %this can be commented out or replaced by a company/department logo

  \includegraphics[height=0.75\headerheight]{preambles/aau_logo_new_neg}
 
 
 %
}
%the poster title
{\color{white}\bf
 \LARGE{
  Day and Night-Time Drive Analysis using Stereo Vision for Naturalistic Driving Studies %Vehicle and Critical Event Detection in Night by Robust Stereo Analysis
 }
}
%the author(s)
{\color{white}\scriptsize
  \vspace{0.7em} \begin{large} \textit{Mark P. Philipsen}, \textit{Morten B. Jensen}, \textit{Andreas Møgelmose}, \textit{Mohan M. Trivedi} \& \textit{Thomas B. Moeslund} \end{large} %\\[0.7em]
   %mphili10@student.aau.dk, mbje10@student.aau.dk\\[0.3em]
%Visual Analysis of People Lab, Aalborg University, Denmark \\
%Computer Vision and Robotics Research Laboratory, University of California, San Diego
%\small\textit{Advisors: Mohan M. Trivedi \& Thomas B. Moeslund}
  %\begin{large}VGIS, Group: 13gr740, Aalborg University, Denmark, 2013\end{large}
}
%the logo (the logo on the right)
{
  %this can be commented out or replaced by a company/department logo
  \includegraphics[height=0.75\headerheight]{preambles/all.png}
%\setlength{\columnsep}{0pt}
%  \begin{tabular}{cc}
%% \subfloat[]{\includegraphics[width=0.46\textwidth]{text/figures/numOfObjects.png} }\\
%% \subfloat[]{\includegraphics[width=0.46\textwidth]{text/figures/numOfObjects.png}}
%%    
%
% \multicolumn{2}{c}{ \includegraphics[height=0.25\headerheight]{preambles/ucsdJacobs} } \\
% & \\
%  \includegraphics[height=0.12\headerheight]{preambles/cvrr} & \includegraphics[height=0.12\headerheight]{preambles/lisaOwn.png} \\
%
%%   \\[1ex]
%
%  
%  \end{tabular}



}
\fontsize{8pt}{8pt}\selectfont

%%%%%%%%%%%%%%%%%%%%%%%%%%%%%%%%%%%%%%%%%%%%%%%%
% the actual content of the poster begins here
%%%%%%%%%%%%%%%%%%%%%%%%%%%%%%%%%%%%%%%%%%%%%%%%

\begin{posterbox}[name=abstract,column=0]{Introduction}
%In order to understand dangerous situations in the driving environment, Naturalistic driving studies (NDS) are conducted by collecting and analyzing data from sensors looking inside and outside of the car. The process of manually analyzing the overwhelming amount of data that is generated in such studies is very comprehensive. This paper proposes a method for data reduction in NDS based on stereo vision during day- and nighttime. The developed system can automatically register six NDS events mainly related to intersections. Additionally, this paper proposes a new drive event relying on the extra dimension provided by stereo vision, where the minimum distance from the ego vehicle to other vehicles directly in front is detected.
In total six drive events are selected on the basis of them being problematic to detect automatically using conventional monocular computer vision. 
Finally the proposed system is evaluated on day- and nighttime data, and the result is a quantified drive analysis report based on the six events. The proposed system have an overall precision of 0.9375 and an overall recall of 0.71, which is considered a useful for data reduction in NDS.

%In order to prevent road accidents, a lot of research has been done in Advanced Driver Assistant Systems (ADAS), where the purpose is to utilize e.g. computer vision systems to assist the driver and prevent accidents while on the road. Most of these systems are, however, evaluated on daytime scenarios, which in 2010 represented 46.2 \% of all accidents on motorways in 20 European countries \cite{euTrafficStats}.
%It is essential to address the nighttime scenarios, as 36.1 \% of motorway accidents happened in low light conditions.

We introduce a stereo-vision system for automatic Naturalistic Driving Study (NDS) data reduction on both day and nighttime data. The purpose of NDS is to provides insight in the patterns and behaviors of drivers during near-crashes and crashes. Since NDS in conducted on large amounts of data, it is extremely time consuming to analyses. Automated solutions, that can sort through and categorize the collected data, are therefore sought after.

A stereo vision based system enables detection and tracking of vehicles in scenarios that would be problematic for monocular detectors. The system proposed in this paper will be limited to handling a handful of NDS events, which especially benefit from the extra dimension in stereo-vision. The NDS events considered in this paper are illustrated in Figure \ref{fig:introduction:semantics}. Monocular systems usually have problems dealing with occlusions and are in some cases using classifiers based on appearance, which require a large amount of training data for dealing with detection of car from the many possible angles and light conditions. 

\begin{figure}[H]
  \centering
  \begin{tabular}{ccc}
    \subfloat[Average number of cars in front of ego-vehicle.]{\includegraphics[width=0.28\textwidth]{text/figures/numOfObjects.png} \label{fig:introduction:semantics:a}} &
    \subfloat[Distance to rear-end of vehicle directly in front.]{\includegraphics[width=0.28\textwidth]{text/figures/avgDistanceSameLane.png} \label{fig:introduction:semantics:b}} &
    
    \subfloat[Other vehicle entering intersection - left turn across path.]{\includegraphics[width=0.28\textwidth]{text/figures/leftIntersect.png} \label{fig:introduction:semantics:c}} \\
    \subfloat[Other vehicle entering intersection - turning onto opposite direction.]{\includegraphics[width=0.28\textwidth]{text/figures/turningO.png} \label{fig:introduction:semantics:d}}&
    
    \subfloat[Other vehicle entering intersection - straight across path.]{\includegraphics[width=0.28\textwidth]{text/figures/passingIntersect.png} \label{fig:introduction:semantics:e}} &
    \subfloat[Other vehicle entering intersection - turning same direction.]{\includegraphics[width=0.28\textwidth]{text/figures/passingIntersectOntoSameDirection.png} \label{fig:introduction:semantics:f}}
  \end{tabular}
  \caption{Drive events that can be automatically detected by the method proposed in this paper. Green car is ego vehicle. Red cars are other vehicles. 
%(a) Average number of cars in front of ego-vehicle (b) Distance to rear-end of vehicle directly in front. (c) Other vehicle entering intersection - left turn across path (d) Other vehicle entering intersection - turning onto opposite direction (e) Other vehicle entering intersection - straight across path (f) Other vehicle entering intersection - turning same direction.
}
\label{fig:introduction:semantics}
\end{figure}

Our contributions are:\\
\begin{itemize}
\item Using stereo-vision for automatic data reduction for NDS on both day and nighttime data, with focus on intersections (Figure \ref{fig:introduction:semantics:c}, \ref{fig:introduction:semantics:d}, \ref{fig:introduction:semantics:e}, \ref{fig:introduction:semantics:f}).
\item Introducing a new NDS event: Average number of vehicles in front of the ego vehicle. (Figure \ref{fig:introduction:semantics:a}).
\item Introducing a new NDS event: Average distance to vehicles directly in front of the ego vehicle. (Figure \ref{fig:introduction:semantics:b}).
\end{itemize}

\vspace{21.5pt}
\end{posterbox}

\begin{posterbox}[name=abstract,column=0]{Related Work}
%In order to understand dangerous situations in the driving environment, Naturalistic driving studies (NDS) are conducted by collecting and analyzing data from sensors looking inside and outside of the car. The process of manually analyzing the overwhelming amount of data that is generated in such studies is very comprehensive. This paper proposes a method for data reduction in NDS based on stereo vision during day- and nighttime. The developed system can automatically register six NDS events mainly related to intersections. Additionally, this paper proposes a new drive event relying on the extra dimension provided by stereo vision, where the minimum distance from the ego vehicle to other vehicles directly in front is detected.
In total six drive events are selected on the basis of them being problematic to detect automatically using conventional monocular computer vision. 
Finally the proposed system is evaluated on day- and nighttime data, and the result is a quantified drive analysis report based on the six events. The proposed system have an overall precision of 0.9375 and an overall recall of 0.71, which is considered a useful for data reduction in NDS.

%We look at some resent research published with regards to automatic NDS data reduction as well as some notable work in improving disparity maps and object detection in traffic scenes. 

%In \cite{Guo20133}, statistics from near-crashes and crashes are used to identify factors related to driver associated risks, such as age, experience, gender, and demographic. Similar studies are conducted in \cite{SimonsMorton2011587}, where results show that near-crash and crash percentage among teenagers were 75 \% lower in the presence of adult passengers and 96 \% higher among teenagers with risky friends. In \cite{mcevoy2006impact} NDS are used to quantify distracting activities from e.g. a mobile phones, which result in loss of concentration. Results show that drivers in average are engaged in distracting activity every 6 minutes, something which could result in a near-crash or crash.



In \cite{raviLaneSemantics}, monocular computer vision and information from the CAN bus is used to automatically detect 23 drive events, including lane position, vehicle localization within lanes, vehicle speed, traffic density, and road curvature. In \cite{raviOvertaking}, a system is developed for automatic labeling of driver behavior events relying on overtaking and receding vehicle detection. This type of system is categorized as Looking-In and Looking-Out (LiLo), which is discussed in depth in \cite{trivedi2007looking}. LiLo fits well with using multiple inputs to understand the driver's behavior. \\%An example of Looking-In is \cite{ohn2013driver} where driver behavior with respect to hand activity is evaluated.\\

%In \cite{SivaramanSurvey}, a review of the research conducted since 2005 in both monocular and stereo vision with respect to vision-based vehicle detection is presented. 
%Before 2005, \cite{bebisReview} conducted a review of on-road vehicle detection. Most of the research published with regard to vehicle detection is evaluating the methods on data, representing a limited part of the challenges. 

\cite{HirschmullerRLandSGBM} introduced the accurate and efficient Semi-Global Matching (SGM) algorithm from finding stereo correspondence. SGM make use of epipolar geometry, and in most cases a set of rectified stereo images. Horizontal lines in the images are used as a scan lines, matches are then found for all possible disparities in a 1D disparity search. A match for a pixel in the left image is found in the right image by searching through the corresponding horizontal row and locating the most similar block to a reference block around the original pixel in the left image. The offset between these pixels is known as the disparity, which, if correctly matched, is directly related to the distance to the corresponding object. A final match it based on the outcome of a smoothed path cost which is calculated in a number of directions for each disparity. In the same paper, the LR-RL-consistency check is proposed for reducing noise in the calculated disparity image. \\

In \cite{labayrade2003onboard}, the so-called V-disparity is generated and used for separating objects from the ground/road surface. The V-disparity examines the vertical coordinates in a (u,v) image coordinate system and is constructed using a disparity map from, e.g. the SGM algorithm. What is especially histograms are calculated for each row in the disparity. Significant surfaces in the disparity map will then show up as lines in the V-disparity. %\cite{labayrade2003onboard} is evaluated on nighttime conditions. 

%In \cite{daimler6DVision}, 6D-vision is introduced, where features are found in the left monocular image and then located in 3D by using the stereo images. For each feature point, a Kalman filter is used to estimate a 6D vector consisting of the 3D position and 3D motion vector. \cite{daimler6DVision} is able to do ego-motion compensation by identifying the static 6D points. The predicted static world points are then compared to the remaining points to isolate the ego-motion. In \cite{daimlerStixel}, this work is continued and by using 6D-vision, tracked feature points are represented as stixels, which is a vertical areas of equal disparity. The tracked objects are being classified using prior knowledge of vehicle shapes. Alternatively, objects can be classified using clustering in the disparity map, as seen in \cite{broggiTerramax}. In \cite{adverseWeatherDaimlar,exploitedTrafficSceneStatistics}, temporal and scene priors from good conditions are used with the purpose of improving the disparity map in adverse weather conditions, such as night, rain, and snow. Using these priors, the object detection rate improves on a database of 3000 frames including bad weather while reducing the false positive rate.
%\vspace{2pt}
\end{posterbox}

\begin{posterbox}[name=propsedMethod,column=1,span=2]{Proposed System}
\vspace*{-1.5mm}
\begin{multicols}{2}   
A Bumblebee XB3 stereo camera is used to acquire \textbf{stereo images} with frame rate of 16 FPS in a resolution of 1280x960 each. For generating the \textbf{disparity map}, OpenCV's CPU implementation of  SGBM is used, it is based on the stereo correspondence algorithm proposed in \cite{HirschmullerRLandSGBM}. \textbf{Noise} is removed using a modified LR-RL consistency check where consistency conflicts are solved by keeping the lowest disparity value, instead of rejecting the pixel entirely. An additional temporal consistency check and a monocular color check is also introduced.
The \textbf{road surface} is found by searching for the most significant line in the V-disparity \cite{labayrade2002real} using RANSAC. Additionally, the line parameters are filtered using a Kalman filter to smooth out faulty road surface detections. Following this the disparity pixels belonging to the road surface or noise below is removed and what remains is provisionally regarded as objects of interest. These objects are projected to 3D using the camera's properties, which result in a 3D \textbf{point cloud} representations of the segmented objects.
%Using the camera's focal length and baseline, along with the calculated disparity, the actual distance to objects in the camera's view can be found. 
The acquired point clouds are post processed using a band-bass filter to remove near and distant points, downsampled using a voxel grid, and outlier removal. \textbf{Clusters} are found by creating a k-d tree, which organize points according to distance to neighbors. 
%and for each point finding neighbors within a specified radius. 
\end{multicols}

%\vspace*{-2.5mm}
\begin{figure}[H]
  \centering
  \includegraphics[width=1\textwidth]{text/figures/methods/flowChart.png}
  \captionof{figure}{Critical event detection process flow.}
  \label{systemOverview::systemOverview}
\end{figure}

%\vspace*{-7mm}
\begin{multicols}{2}   
The clusters' center points are used for nearest neighbor \textbf{tracking} between frames and for determining the distance from ego-vehicle to detected vehicles.
For \textbf{ego-motion compensating} while the ego-vehicle is approaching an intersection, we utilize \textit{LIBVISO2: C++ Library for Visual Odometry 2} \cite{Geiger2011IV}. 
Events are detected by looking at the \textbf{movement} history of other vehicles. For all detected vehicles, individual frame to frame movements are categorized to form a histogram of movements for the sequence, which is used for \textbf{event classification}.
An overview of the system is shown in Figure \ref{systemOverview::systemOverview}. The final system output is an event report.


\end{multicols}
%\vspace*{1pt}


%\vspace*{2pt}

%
%
%\begin{figure}[H]
%  \centering
%  \includegraphics[width=0.42\textwidth]{text/figures/systemForeslag4.png}
%  \captionof{figure}{Overview of the methods used in the proposed NDS system.}
%  \label{systemOverview::systemOverview}
%\end{figure}
%
%%The road surface is removed using the detected line in a corresponding v-disparity. The remaining pixels are considered vehicle candidate pixels. These pixels are projected into 3D world coordinates and outliers are removed. Segmentation into clusters is done by grouping points that are close neighbors into the same clusters. Each sufficiently big cluster is regarded as a detected vehicle. Finally, the detected vehicles are tracked and NDS event are detected and logged in a drive analysis report. \\[5pt]
%
%\textbf{Generate Disparity Map \& Noise Removal}\\
%A disparity map is generated using OpenCV's SGBM implementation. For noise reduction, a LR-RL consistency check is done \cite{HirschmullerRLandSGBM}. In case of a disparity differences in the consistency check, the lowest value is selected as the output pixel. A temporal consistency check assumes that a pixel's disparity value does not change erratically between frames. Finally, a monocular color check removes pixels from extremely light and dark regions.
%
%\textbf{Road Surface Removal}\\
%The road surface is found by searching for the most significant line in the V-disparity\cite{labayrade2002real} using RANSAC. Additionally, the line parameters are filtered using a Kalman filter to smooth out faulty road surface detections. The parameters are then used to remove pixels for everything not above the road surface. The remaining pixels are considered vehicle candidate disparity pixels. In Figure \ref{fig:methods::disparity:dispEx1:postVdisp} the left input image and the vehicle candidate disparity map are shown.
%
%\begin{figure}[H]
%  \centering
%  \begin{tabular}{cc}
%    \subfloat[A rectified image captured]{\includegraphics[width=0.22\textwidth]{text/figures/methods/dispEx1/inputL.png} \label{fig:methods::disparity:dispEx1:inputl}} &
%    \subfloat[Noise reduced disparity image.]{\includegraphics[width=0.22\textwidth]{text/figures/methods/dispEx1/postVDisp.png} \label{fig:methods::disparity:dispEx1:postVDisp}}\\
%  \end{tabular}
%  \caption{Results after Generate Disparity Map \& Noise Removal and Road Surface Removal.}
%\label{fig:methods::disparity:dispEx1:postVdisp}
%\end{figure}
%
%\textbf{Project To Point Cloud \& Cloud Filtering}\\
%Using the camera's focal length $f$(in pixels) and it's baseline $b$(in meters), along with the calculated disparity $d$(in pixels), the actual distance $z$ to objects in the camera's view can be found. The remaining $x$ and $y$ component of the world coordinate are found using the column and row index of the disparity pixel.% \\[5pt]
%
%\textbf{Segmentation \& Tracking}\\
%The acquired point clouds are post processed using a band-bass filter to remove near and distant points, downsampled using a voxel grid, and removing outliers. Clusters are found by creating k-d tree, and for each point finding neighbors within a specified radius. The clusters center point is used for nearest neighbor tracking between frame and used to determine the distance from ego vehicle to detected vehicle.
%
%
%\textbf{NDS Event Detection}\\
%NDS events are detected by looking at the movement history of other vehicles with regard to the ego vehicle. For all detected vehicles, individual frame to frame movements are categorized to form a histogram of events for determining which NDS events have occurred. 
%
%\begin{figure}[H]
%  \centering
%  \begin{tabular}{cc}
%    \subfloat[]{\includegraphics[width=0.22\textwidth]{text/figures/methods/intersectionVoting.png} \label{fig:methods::intersectionVoting}} &
%    \subfloat[]{\includegraphics[width=0.22\textwidth]{text/figures/methods/frame2.png} \label{fig:methods::disparity:dispEx1:postVDisp}}\\
%  \end{tabular}
%  \caption{(a) Illustration of various movements that are detected by the system. (b) NDS detection example.}
%\label{fig:methods::disparity:dispEx1:postVdisp}
%\end{figure}
%
%For all detected vehicles, individual frame to frame movements are categorized to form a histogram of events for determining which NDS events have occurred. 
%An example of such a histogram is seen in Figure \ref{methods::histogram}.
%\begin{figure}[H]
%\centering
%     \includegraphics[width=0.35\textwidth]{text/figures/methods/eventHistogram.png}
%     \caption{Example of histogram of detected movements.}
%     \label{methods::histogram}
%\end{figure}




\end{posterbox}

\begin{posterbox}[name=results,column=1,below=propsedMethod]{Results}
\begin{multicols}{2}   
The proposed system is evaluated on 4,992 day and 3,933 night time frames. In Table \ref{systemOverview::derivedSemantics:Results} the results are seen, where GT is the \textit{ground truth} of events manually labeled and SO is the \textit{system output}. The proposed system have an overall precision of 0.78 and recall of 0.72.
%\vspace*{-5mm}
\begin{center}
\captionof{table}{Summary of event report analysis. The syntax for the results are [SO/GT]. P and R are abbreviations for precision and recall, respectively.}
\resizebox{0.45\textwidth}{!}{%
  \begin{tabular}{l r r r r r}
   \toprule
 \textbf{\textit{Drive Behavior Event}}& \textbf{Daytime} & \textbf{Nighttime} & \textbf{P} & \textbf{R} \\
    \toprule
    From Left To Right - straight across path & 35/32 & 5/19 & 0.95 & 0.63 \\ 
    From Right to Left - straight across path & 45/34  &  11/33  & 0.87 & 0.67 \\
    Left turn across path  & 5/5 & 20/1  & 0.75 & 1 \\ 
    Turn onto opposite dir.  & 32/37  & 41/15  & 0.68 & 0.93 \\ 
    Short turn onto same dir. & 7/5  & 9/5 & 0.63 & 1 \\ 
    Long turn onto same dir. & 1/16  & 1/8  &  1 & 0.09 \\ 
    Avg. number of cars & 1.67/1.74 & 1.6/1.3 &  NA &NA  \\ 
    Avg. distance to car &  8.73 m &10.98 m &  NA& NA \\ 
    \bottomrule     
    \label{systemOverview::derivedSemantics:Results}   
  \end{tabular}
  
}
\end{center}
In Figure \ref{systemOverview::distanceToVehicle} an example of the output from the system is seen. %This specific image is an example of the overlapping
The new NDS event introduced with this work is the distance to vehicle in front of the ego-vehicle, which is measured using stereo-vision. In Figure \ref{systemOverview::distanceToVehicle} an example of this is seen where the distance of 12.6 meter is written in blue just above the center. This distance measure is also visible in Figure \ref{systemOverview::occludedVehicleDetected}.
\begin{figure}[H]
  \centering
  \includegraphics[width=0.3\textwidth]{text/figures/lefttrun2cars.png}
  \captionof{figure}{Partially occluded vehicle detected in a left turn.}
  \label{systemOverview::occludedVehicleDetected}
\end{figure}

\begin{figure}[H]
  \centering
  \includegraphics[width=0.3\textwidth]{text/figures/distanceCarEvent.png}
  \captionof{figure}{Distance to vehicle in front of ego-vehicle.}
  \label{systemOverview::distanceToVehicle}
\end{figure}
\end{multicols}
\end{posterbox}

\begin{posterbox}[name=conclusion,column=2,below=propsedMethod]{Concluding Remarks}
The use of stereo-vision is considered beneficial, especially in scenarios with partly occluded cars, in such cases most monocular systems will fail. It is considered that providing distance information from the ego vehicle to surroundings provide useful information with regard to drive patterns before and during a crash or near crash. The system proved to work in both day and nighttime conditions with a limited drop in overall performance for the night scenarios. The proposed system have an overall precision of 0.9375 and an overall recall of 0.71, which is considered useful for data reduction in NDS. Future work may include looking at additional NDS events where stereo-vision can be utilized.

%\vspace{8pt}
\end{posterbox}

\begin{posterbox}[name=acknowledgment,column=2,below=conclusion]{Acknowledgment}
Thanks to Ravi K. Satzoda and Andreas Møgelmose for guidance.
%\vspace{5pt}
\end{posterbox}

\let\oldbibliography\thebibliography
\renewcommand{\thebibliography}[1]{\oldbibliography{#1}
\setlength{\itemsep}{0pt}} %Reducing spacing in the bibliography.
\begin{posterbox}[name=references,column=0,below=abstract]{References}
%%%In order to prevent road accidents, a lot of research has been done in Advanced Driver Assistant Systems (ADAS), where the purpose is to utilize e.g. computer vision systems to assist the driver and prevent accidents while on the road. Most of these systems are, however, evaluated on daytime scenarios, which in 2010 represented 46.2 \% of all accidents on motorways in 20 European countries \cite{euTrafficStats}.
%It is essential to address the nighttime scenarios, as 36.1 \% of motorway accidents happened in low light conditions.

We introduce a stereo-vision system for automatic Naturalistic Driving Study (NDS) data reduction on both day and nighttime data. The purpose of NDS is to provides insight in the patterns and behaviors of drivers during near-crashes and crashes. Since NDS in conducted on large amounts of data, it is extremely time consuming to analyses. Automated solutions, that can sort through and categorize the collected data, are therefore sought after.

A stereo vision based system enables detection and tracking of vehicles in scenarios that would be problematic for monocular detectors. The system proposed in this paper will be limited to handling a handful of NDS events, which especially benefit from the extra dimension in stereo-vision. The NDS events considered in this paper are illustrated in Figure \ref{fig:introduction:semantics}. Monocular systems usually have problems dealing with occlusions and are in some cases using classifiers based on appearance, which require a large amount of training data for dealing with detection of car from the many possible angles and light conditions. 

\begin{figure}[H]
  \centering
  \begin{tabular}{ccc}
    \subfloat[Average number of cars in front of ego-vehicle.]{\includegraphics[width=0.28\textwidth]{text/figures/numOfObjects.png} \label{fig:introduction:semantics:a}} &
    \subfloat[Distance to rear-end of vehicle directly in front.]{\includegraphics[width=0.28\textwidth]{text/figures/avgDistanceSameLane.png} \label{fig:introduction:semantics:b}} &
    
    \subfloat[Other vehicle entering intersection - left turn across path.]{\includegraphics[width=0.28\textwidth]{text/figures/leftIntersect.png} \label{fig:introduction:semantics:c}} \\
    \subfloat[Other vehicle entering intersection - turning onto opposite direction.]{\includegraphics[width=0.28\textwidth]{text/figures/turningO.png} \label{fig:introduction:semantics:d}}&
    
    \subfloat[Other vehicle entering intersection - straight across path.]{\includegraphics[width=0.28\textwidth]{text/figures/passingIntersect.png} \label{fig:introduction:semantics:e}} &
    \subfloat[Other vehicle entering intersection - turning same direction.]{\includegraphics[width=0.28\textwidth]{text/figures/passingIntersectOntoSameDirection.png} \label{fig:introduction:semantics:f}}
  \end{tabular}
  \caption{Drive events that can be automatically detected by the method proposed in this paper. Green car is ego vehicle. Red cars are other vehicles. 
%(a) Average number of cars in front of ego-vehicle (b) Distance to rear-end of vehicle directly in front. (c) Other vehicle entering intersection - left turn across path (d) Other vehicle entering intersection - turning onto opposite direction (e) Other vehicle entering intersection - straight across path (f) Other vehicle entering intersection - turning same direction.
}
\label{fig:introduction:semantics}
\end{figure}

Our contributions are:\\
\begin{itemize}
\item Using stereo-vision for automatic data reduction for NDS on both day and nighttime data, with focus on intersections (Figure \ref{fig:introduction:semantics:c}, \ref{fig:introduction:semantics:d}, \ref{fig:introduction:semantics:e}, \ref{fig:introduction:semantics:f}).
\item Introducing a new NDS event: Average number of vehicles in front of the ego vehicle. (Figure \ref{fig:introduction:semantics:a}).
\item Introducing a new NDS event: Average distance to vehicles directly in front of the ego vehicle. (Figure \ref{fig:introduction:semantics:b}).
\end{itemize}

\vspace{21.5pt}
\nocite{bal:cha:gra:pae}
\bibliographystyle{IEEEtran}
\scriptsize{
\bibliography{literature}}
\end{posterbox}

\end{poster}

\end{document}
